\pdfoutput=1
\documentclass[11pt]{article}
\usepackage{smile}
\usepackage{subcaption}
\usepackage{pdflscape}   
\usepackage{pifont}
\usepackage{xcolor}
\usepackage{tikz}
\usetikzlibrary{arrows.meta, positioning, calc, shapes.geometric}
% ===== Morandi palette =====
\definecolor{morandiGreen}{RGB}{140,160,140}
\definecolor{morandiPink}{RGB}{244,186,180}  
\definecolor{morandiBlue}{RGB}{190,210,250}  
\definecolor{morandiOrange}{RGB}{200,160,120}
\definecolor{morandiPurple}{RGB}{150,140,160}
\definecolor{morandiGray}{RGB}{160,160,160}
\usepackage[margin=1in]{geometry}
\usepackage[most]{tcolorbox}
\renewcommand{\baselinestretch}{1.10}
%%%%%%%%%%%%-- Comments --%%%%%%%%%%%%
%%%%%%%%%%%%-- new command --%%%%%%%%%%%%
\newcommand{\bbP}{\mathbb{P}}
\newcommand{\bbR}{\mathbb{R}}
\newcommand{\bbE}{\mathbb{E}}
\newcommand{\bbN}{\mathbb{N}}
\newcommand{\wh}{\widehat}
\newcommand{\wt}{\widetilde}
\newcommand{\ol}{\overline}
%%%%%%%%%%%%-- start --%%%%%%%%%%%%%%%
\title{\bf GAN and Markov Chain Monte Carlo}
%\author{Junyi Liao$^\dagger$
%\\ 
%$^{\dagger}$\small\it Department of Electrical and Computer Engineering, Duke University
%}
% \author{Junyi Liao$^1$,\ \ 
% Xiaohui Jiang$^2$
% \\ 
% ~\\
% $^{1}$\small\it Department of Electrical and Computer Engineering, Duke University \\ 
% $^{2}$\small\it Department of Biostatistics and Bioinformatics, Duke University \\ 
% \small\sf \{junyi.liao, x.jiang\}@duke.edu
% }
\date{}
\begin{document}
\maketitle

%\begin{abstract}
%\end{abstract}

%\tableofcontents

\section{Motivation and Problem Statement}
Let $\pi$ be a target probability distribution on a state space $\mathsf X$. Our main setting assumes that we can obtain i.i.d. training samples from $\pi$ but cannot evaluate its density. We seek an efficient, reusable Markov sampler that can generate new samples from arbitrary initial states. In a secondary setting, the target may be specified through an unnormalized energy $\pi(x)\propto\exp\{-U(x)\}$.

A classical GAN learns a static mapping $G_\theta(Z)$ whose marginal distribution approximates $\pi$. By contrast, we consider a \emph{state-dependent} generator
\begin{equation}
Y_t=G_\theta(X_t,Z_t),\qquad Z_t\sim p_Z,
\label{eq:proposal}
\end{equation}
which induces a possibly implicit proposal kernel $q_\theta(dy\mid x)$. A separate neural function $a_\psi(x,y)\in[0,1]$ determines whether the proposal is accepted:
\begin{equation}
X_{t+1}=\begin{cases}
Y_t, & U_t\le a_\psi(X_t,Y_t),\\
X_t, & \text{otherwise},
\end{cases}\qquad U_t\sim\operatorname{Unif}(0,1).
\label{eq:chain}
\end{equation}
Unlike a standard GAN discriminator, $a_\psi$ assesses a \emph{transition} $(x,y)$, not just an isolated sample. The proposal and acceptance mechanisms are naturally cooperative: they should produce a chain that both preserves $\pi$ and explores it rapidly. An adversarial critic can be used to identify failures of these objectives.

\paragraph{Goal.} Develop a sample-based minimax formulation that jointly learns $(q_\theta,a_\psi)$ and controls (i) the target-distribution error, (ii) violations of stationarity or detailed balance, and (iii) sampling inefficiency such as autocorrelation or mode trapping.

\section{Relevant Prior Work}
\paragraph{MH-GAN: independent proposals and discriminator correction.}
\cite{turner2019metropolis} treat a pre-trained GAN generator $p_\theta$ as an \emph{independent} proposal. For a fixed generator, the Bayes-optimal GAN discriminator satisfies
\begin{equation}
D^*(x)=\frac{\pi(x)}{\pi(x)+p_\theta(x)},
\qquad
\frac{\pi(x)}{p_\theta(x)}=\frac{D^*(x)}{1-D^*(x)}.
\end{equation}
Hence the independent MH acceptance rule can be written without evaluating either density:
\begin{equation}
a_{\mathrm{MH}}(x,y)
=\min\left\{1,\frac{D^*(y)/[1-D^*(y)]}{D^*(x)/[1-D^*(x)]}\right\}.
\label{eq:mhgan}
\end{equation}
This is an MH \emph{wrapper} around a generator, rather than a jointly trained Markovian proposal and acceptance rule. Exactness requires the correct density ratio (and appropriate support and convergence assumptions). The paper also stresses discriminator calibration, which matters because a good classification ranking need not produce an accurate MH ratio.

\paragraph{Implicit MH: state-dependent implicit proposals.}
\cite{neklyudov2019implicit} generalize the idea to a proposal $q_\theta(y\mid x)$ that need not have a tractable density. Introduce the forward and reverse joint proposal measures (written as densities here for convenience):
\begin{equation}
F_\theta(x,y)=\pi(x)q_\theta(y\mid x),
\qquad
B_\theta(x,y)=\pi(y)q_\theta(x\mid y)=F_\theta(y,x).
\label{eq:flows}
\end{equation}
A pairwise discriminator that classifies samples from $F_\theta$ versus $B_\theta$ has optimum
\begin{equation}
D^*_{\theta}(x,y)=\frac{F_\theta(x,y)}{F_\theta(x,y)+B_\theta(x,y)},
\end{equation}
so the desired acceptance ratio becomes
\begin{equation}
a_{\mathrm{MH}}(x,y)=\min\left\{1,\frac{B_\theta(x,y)}{F_\theta(x,y)}\right\}
=\min\left\{1,\frac{1-D^*_{\theta}(x,y)}{D^*_{\theta}(x,y)}\right\}.
\label{eq:implicit-mh}
\end{equation}
Crucially, one can sample a forward pair $(X,Y)$ by taking $X\sim\pi$ and $Y\sim q_\theta(\cdot\mid X)$; the reversed pair $(Y,X)$ follows $B_\theta$. Thus neither density needs to be evaluated. The paper analyzes the stationary-distribution error caused by an imperfect discriminator, under explicit mixing/minorization conditions, and considers both independent and Markov proposals. 

\paragraph{Adversarially learned Markov kernels.}
\cite{song2017generative} learn a Markov transition through adversarial comparisons with data, including short chains initialized from target samples. Their subsequent A-NICE-MC approach \cite{song2017nice} learns neural MCMC proposals and uses an analytically specified MH correction for energy-based targets. \cite{neklyudov2018metropolis} also connect GANs, variational inference, and MH acceptance objectives. Together, these works already cover adversarial Markov generators, learned proposals, and discriminator-based MH corrections. A joint proposal--acceptance game must be distinguished from these precedents through its specific objective or guarantees.

\section{Proposed Proposal--Acceptance--Critic Framework}
\subsection{Neural transition kernel}
Given \eqref{eq:proposal}--\eqref{eq:chain}, the resulting Markov kernel is
\begin{equation}
K_{\theta,\psi}(x,dy)=q_\theta(dy\mid x)a_\psi(x,y)
+\left[1-\int q_\theta(dz\mid x)a_\psi(x,z)\right]\delta_x(dy).
\label{eq:kernel}
\end{equation}
We propose three interacting components: (i) $G_\theta$ generates state-dependent candidate moves, (ii) $a_\psi$ learns a pair-dependent accept/reject probability, and (iii) a multi-head critic evaluates distributional mismatch, accepted-flow asymmetry, and slow mixing. Only the \emph{critic} is adversarial to the sampler; proposal and acceptance jointly minimize the sampler's objective.

\subsection{Matching the target after a finite number of steps}
Let $\nu_0$ be an initialization distribution and $\mu_{\theta,\psi}^{(T)}=\nu_0K_{\theta,\psi}^{T}$. An adversarial integral probability metric is
\begin{equation}
\mathcal L_{\mathrm{gen}}(\theta,\psi)=\sup_{f\in\mathcal F}
\left\{\bbE_{X\sim\pi}f(X)-\bbE_{X\sim\mu_{\theta,\psi}^{(T)}}f(X)\right\}.
\label{eq:gen}
\end{equation}
For instance, $\mathcal F$ may be a bounded-Lipschitz class; with $1$-Lipschitz critics and finite first moments, the objective is Wasserstein-1. A bounded class is especially relevant when targeting heavy-tailed distributions that may lack a finite first moment. Matching a single finite-time marginal is not, by itself, a stationarity guarantee.

\subsection{Adversarial detailed-balance constraint}
Define the \emph{accepted flow}
\begin{equation}
J_{\theta,\psi}(dx,dy)=\pi(dx)q_\theta(dy\mid x)a_\psi(x,y).
\label{eq:accepted-flow}
\end{equation}
Detailed balance is exactly the symmetry condition $J_{\theta,\psi}(dx,dy)=J_{\theta,\psi}(dy,dx)$. We therefore propose the sample-estimable discrepancy
\begin{align}
\mathcal L_{\mathrm{bal}}(\theta,\psi)
&=\sup_{h\in\mathcal H}\int\!\big[h(x,y)-h(y,x)\big]J_{\theta,\psi}(dx,dy)\notag\\
&=\sup_{h\in\mathcal H}\bbE_{X\sim\pi,\,Y\sim q_\theta(\cdot\mid X)}
\Big[a_\psi(X,Y)\{h(X,Y)-h(Y,X)\}\Big].
\label{eq:balance}
\end{align}
Here $\mathcal H$ should be symmetric under sign reversal and suitably bounded. With all bounded measurable test functions, zero discrepancy is equivalent to detailed balance; with a neural critic, only the restricted discrepancy is controlled. The loss requires neither $\pi(x)$ nor $q_\theta(y\mid x)$.

\paragraph{A degeneracy.} The identically rejecting rule $a_\psi\equiv0$ gives $\mathcal L_{\mathrm{bal}}=0$, yet the chain never moves. We must therefore explicitly reward exploration, rather than use detailed balance as the sole learning criterion.

\subsection{An adversarial mixing-efficiency objective}
When $K_{\theta,\psi}$ is $\pi$-reversible, its Poincar\'e (right) spectral gap admits the Dirichlet-form representation
\begin{align}
\operatorname{Gap}(K_{\theta,\psi})
&=\inf_{\substack{\bbE_\pi g=0\\\operatorname{Var}_\pi(g)=1}}
\mathcal E_{K_{\theta,\psi}}(g,g),\\
\mathcal E_{K_{\theta,\psi}}(g,g)
&=\frac12 \bbE_{X\sim\pi,\,Y\sim K_{\theta,\psi}(\cdot\mid X)}
\big[(g(Y)-g(X))^2\big]\\
&=\frac12 \bbE_{X\sim\pi,\,Y\sim q_\theta(\cdot\mid X)}
\Big[a_\psi(X,Y)(g(Y)-g(X))^2\Big].
\label{eq:dirichlet}
\end{align}
The last expression avoids simulating the discrete accept/reject decision and is differentiable in the acceptance score (and in a reparameterizable proposal). We can approximate the infimum with normalized neural test functions $g\in\mathcal G$. The sampler tries to \emph{maximize} the worst-case Dirichlet energy, while the critic seeks slowly evolving observables. For reliable convergence bounds one may additionally impose a lazy kernel or another condition controlling negative eigenvalues. When detailed balance holds only approximately, this should be treated as a mixing surrogate rather than an exact spectral gap.

\subsection{A candidate minimax objective}
A unified, \emph{provisional} two-player game is
\begin{align}
\min_{\theta,\psi}\max_{\substack{f\in\mathcal F,\ h\in\mathcal H,\\
                                     g\in\mathcal G:\ \bbE_\pi g=0,\ \operatorname{Var}_\pi(g)=1}}
\quad &\underbrace{\bbE_\pi f-\bbE_{\nu_0K_{\theta,\psi}^{T}}f}_{\text{finite-time distribution match}}
+\lambda_{\mathrm{bal}}
\underbrace{\bbE_{\pi q_\theta}\!\left[a_\psi(X,Y)(h(X,Y)-h(Y,X))\right]}_{\text{accepted-flow balance}}
\notag\\[-0.2em]
&-\lambda_{\mathrm{mix}}
\underbrace{\frac12\bbE_{\pi q_\theta}\!\left[a_\psi(X,Y)(g(Y)-g(X))^2\right]}_{\text{worst-case mixing surrogate}}.
\label{eq:game}
\end{align}
The critic may be implemented as a neural network with three heads $(f,h,g)$, or as separate networks. $\lambda_{\mathrm{bal}},\lambda_{\mathrm{mix}}>0$ express competing priorities. Because a finite penalty does not enforce exact balance, a complementary formulation is to maximize the mixing criterion \emph{subject to} a prescribed bound $\mathcal L_{\mathrm{bal}}\le\varepsilon$, and then study the effect of $\varepsilon$ on stationary-distribution error. The objective in \eqref{eq:game} is a research hypothesis, not an established optimization guarantee.

\subsection{Acceptance-network design and possible training strategy}
We propose comparing two variants:
\begin{enumerate}
\item \textbf{Density-ratio-informed acceptance:} learn a pairwise classifier $D_\eta$ from forward/reversed proposals and set $a(x,y)=\min\{1,(1-D_\eta(x,y))/D_\eta(x,y)\}$, optionally with calibration and balance/mixing regularization. This is closely related to implicit MH, but one can investigate \emph{joint} rather than sequential training of the proposal.
\item \textbf{Directly learned acceptance:} parameterize $a_\psi(x,y)=\sigma(s_\psi(x,y))$ and train it through \eqref{eq:game}, with explicit reversibility constraints or penalties. This gives a more flexible objective, but exact stationarity is no longer automatic.
\end{enumerate}
A practical first stage can train the acceptance classifier using real starting states and forward/reversed pairs. A second stage can update proposals and acceptance jointly using \eqref{eq:balance} and \eqref{eq:dirichlet}, adding short-chain endpoint matching once stable gradients are available. Differentiating through discrete multi-step accept/reject decisions remains a technical challenge (e.g., score-function estimators or carefully analyzed relaxations).

\section{What Guarantees Are Plausible?}
\paragraph{Exact balance implies stationarity.}
If the accepted-flow symmetry in \eqref{eq:accepted-flow} holds exactly, then $\pi K_{\theta,\psi}=\pi$, independently of whether the proposal density is evaluable. Ergodicity and convergence from arbitrary initial states require additional irreducibility/recurrence and aperiodicity or related assumptions.

\paragraph{Approximate balance and stationary error.}
Suppose the learned kernel has a unique invariant distribution $\pi_{\theta,\psi}$ and is contractive in total variation with coefficient $\rho<1$. Then
\begin{equation}
\|\pi_{\theta,\psi}-\pi\|_{\TV}
\le\frac{\|\pi K_{\theta,\psi}-\pi\|_{\TV}}{1-\rho}.
\label{eq:stability}
\end{equation}
The numerator is controlled by a sufficiently strong accepted-flow asymmetry measure; it need not be controlled by a finite-capacity neural critic without additional approximation assumptions. This suggests studying a three-way tradeoff between discriminator capacity, balance error, and mixing/contraction.

\paragraph{Why learn acceptance at all?}
For a \emph{fixed} proposal and exact detailed balance, every valid acceptance probability satisfies
\begin{equation}
a(x,y)\le\min\left\{1,\frac{\pi(y)q(x\mid y)}{\pi(x)q(y\mid x)}\right\}=a_{\mathrm{MH}}(x,y)
\end{equation}
where the ratio is defined. Thus classical MH is already pointwise maximal in this class. Learning acceptance is principally motivated by \emph{inaccessible density ratios}, not by the prospect of improving upon exact MH at a fixed proposal. Efficiency improvements must primarily arise from learning better proposals, considering computational costs, or moving beyond the ordinary reversible accept/reject class.

\section{Key Research Questions and Initial Experiments}
\begin{enumerate}[label=\textbf{Q\arabic*.}]
\item \textbf{Novelty and identifiability.} What objective or guarantee distinguishes jointly trained $(q_\theta,a_\psi)$ from MH-GAN, implicit MH, adversarial Markov-chain learning, and A-NICE-MC? Is the acceptance function identifiable, or only its induced kernel?
\item \textbf{Correctness.} Under what function-class, support, and ergodicity assumptions does a small adversarial balance loss imply a small stationary-distribution error? Can a meaningful finite-sample bound include both critic estimation error and proposal-learning error?
\item \textbf{Efficiency.} Can a neural Dirichlet-form objective increase spectral gap or effective sample size per unit computation without sacrificing stationarity? Which observables should the neural mixing critic cover, especially for rare modes and heavy tails?
\item \textbf{Optimization.} Can training avoid the zero-acceptance solution, poor pairwise-classifier calibration, vanishing gradients, and excessive variance from discrete accept/reject steps? Do we need alternating, bilevel, or staged optimization?
\item \textbf{Scope.} Should the main theory assume a fixed trained kernel, or extend to online adaptation? Can the framework handle heavy-tailed and multimodal targets, or nonreversible/augmented-state kernels?
\end{enumerate}

\paragraph{First experimental program.}
Start with low-dimensional multimodal Gaussian mixtures and Student-$t$ mixtures, where ground-truth diagnostics are available. Compare (i) independent GAN sampling, (ii) MH-GAN, (iii) implicit MH with a fixed Markov proposal, (iv) adversarially learned Markov kernels, and (v) joint proposal--acceptance training, with exact MH as a reference whenever target/proposal densities are available. Report target-distribution mismatch, mass of individual modes, acceptance rate, autocorrelation or function-specific ESS, transition cost, and sensitivity to critic capacity and calibration. Distinguish \emph{stationary} errors from errors induced by finite burn-in or finite-length sampling.

\paragraph{Immediate decision for the next meeting.}
We suggest first choosing between two clearly different goals: 
\begin{itemize}[itemsep=0pt,topsep=0pt]
\item[(A)] jointly optimize a proposal with a classifier-based approximation of the standard MH rule, or 
\item[(B)] directly learn a more general acceptance function under adversarial stationarity constraints.
\end{itemize} 
Goal (A) has a clearer interpretation and stronger prior results; goal (B) is broader but entails a substantially harder correctness problem. Both should be evaluated against the close precedents above before making a novelty claim.


%%%%%%%%%%%%-- main content --%%%%%%%%%%
%\newpage




%%%%%%%%%%%%-- reference --%%%%%%%%%%%
\bibliographystyle{ims}
\bibliography{reference}

\newpage
\appendix
\end{document}

 
